\section{Introduction}

\subsection{Background: VED and IFGT as theoretical infrastructure}

This paper is a self-contained structural argument, but its descriptive vocabulary and dynamical syntax are inherited from two prior frameworks.

The first is \textbf{Vortical Enclosure Dynamics (VED)} \citep{VEDVol1,VEDVol2,VEDVol3,VEDVol4}. VED begins from the single foundational axiom that there is difference. It describes time, space, matter, gravity, and cosmological structure as consequences of processes that tend toward closure without becoming complete closure. In the present paper, what is used from VED is its \textbf{differential-development structure}: a non-closed structure continues to transform under its own dynamics.

The second is \textbf{Information Field Geometry Theory (IFGT)} \citep{IFGT}. IFGT provides a minimal vocabulary of closure degree $C$, information density $I$, information flow $J$, and information-field potential $\Phi$. This paper uses IFGT not as a completed explanatory theory, but as a descriptive vocabulary that does not import normativity, intentionality, or subjectivity into the analysis. Here $I$, $J$, and $\Phi$ should be read as projected variables within the intelligence-layer observational window, not as direct identifications with the general IFGT fields.

There is an important methodological reason for using these two frameworks. Existing vocabularies of intelligence often import concepts such as ability, purpose, meaning, judgment, correctness, or agency. These concepts may be indispensable at later interpretive layers, but they are too strong for the structural level considered here. VED and IFGT provide a lower-level language in which motion, accumulation, constraint, and non-closure can be described without prematurely closing the interpretation.

The central claim of this paper, namely the three minimal axioms introduced in \S{}3, does not logically depend on VED or IFGT. The axioms could in principle be written in other dynamical vocabularies. VED and IFGT are adopted because they provide the least obstructive medium currently available for articulating these constraints without importing normativity.

Several parallel papers in the series are relevant. \textbf{Bio-IFGT} \citep{BioIFGT} treats quasi-closure regimes at the biological layer and supplies the biological vocabulary used in the four-layer correspondence in \S{}9.4. The \textbf{consciousness paper} \citep{Consciousness} treats fast/slow log referencing within individual temporal non-closure. The local horizon-form of intelligence discussed in this paper and the consciousness window discussed there are different projections of the same differential-development structure. The word ``horizon'' does not imply multiple absolute boundaries. It means that the \DH{} appears locally in different observational windows.

\subsection{The intelligence question: why neither ability nor knowledge suffices}

The question ``what is intelligence?'' has traditionally been answered in two broad ways. One identifies intelligence with an ability. The other identifies intelligence with knowledge. Both approaches run into a common structural difficulty.

The difficulty of the ability view is that the object of the ability must be specified. To say that intelligence is the ability to do something already requires a criterion for what should count as the relevant thing to do. Problem solving, adaptation, abstraction, planning, and generalization are all candidates, but the selection of any candidate already depends on a prior standard of intelligence.

The difficulty of the knowledge view is familiar from the Gettier problem \citep{Gettier1963}. The classical definition of knowledge as justified true belief produces indefinitely many counterexamples. Adding further stopping conditions does not end the problem; it generates new points at which the definition can fail.

These difficulties are not unrelated. They share the same structural defect: both treat intelligence as a \textbf{stopping structure}. Ability is treated as something one has or lacks. Knowledge is treated as something that is established or not established. In both cases, intelligence is made into a static possession.

This paper reverses that assumption:

\begin{quote}
Intelligence is not a stopping structure. It is the motion that continually produces stopping structures.
\end{quote}

The reversal is not merely terminological. If stopping-structure accounts of intelligence repeatedly destabilize, then the shift to a dynamical account is a structural requirement.

This position has also been reached independently in biophysics. Research on \emph{Physarum polycephalum} since Nakagaki et al. (2000) has described intelligent behavior without attributing intelligence as a possessed property \citep{Nakagaki2000,Tero2007}. Nakagaki's formulation, that slime mold does not possess intelligence but behaves intelligently, converges with the present paper's central position from a very different research path.

\subsection{Contributions and scope}

The main contributions of this paper are:

\begin{enumerate}
\item \textbf{Layer separation} (\S{}2): perception, recognition, knowledge, and intelligence are separated as distinct structural layers.
\item \textbf{Three minimal axioms} (\S{}3): connected state space, temporal ordering, and global constraint $\Phi$ are formulated as simultaneous necessary conditions for intelligence-like dynamics.
\item \textbf{Vocabulary foundation} (\S{}4): IFGT is positioned as a non-normative descriptive vocabulary and related lexically to the Tero--Kobayashi--Nakagaki Physarum solver.
\item \textbf{Cross-domain structural comparison} (\S{}5): slime mold and Transformer dynamics are compared as structurally corresponding systems under the three axioms.
\item \textbf{Reinterpretation of the Gettier problem} (\S{}6): the Gettier problem is read as a diagnostic of mismatch between stopping structures and ongoing dynamics.
\item \textbf{Instability of accumulated correctness} (\S{}7): the accumulation of correct knowledge is shown to generate branching explosion.
\item \textbf{Individual non-closure} (\S{}8): recursive causal inference is used to argue that no stopping operator exists inside the individual system.
\item \textbf{Naming of the boundary} (\S{}9): the local horizon-form called the differential horizon of intelligence is defined and placed in a four-layer correspondence.
\end{enumerate}

The paper does not address the following:

\begin{itemize}
\item social institutions, collective decision-making, or governance, except as future domains in which externalized stopping structures must be constructed;
\item ethical or normative prescriptions;
\item implementation protocols or performance metrics;
\item consciousness, subjectivity, or experiential content;
\item the concrete social and institutional structures that appear after externalization.
\end{itemize}

The paper formulates constraints rather than final answers. The six constraints summarized in \S{}10 are necessary conditions. A theory of intelligence that fails to satisfy them will tend to reproduce the instabilities described here.
